\chapter{Введение}

Современные веб-приложения, предназначенные для работы с медиаконтентом, как правило, реализуются по клиент-серверной архитектуре. Пользователь взаимодействует с интерфейсом через браузер, который отправляет запросы к серверной части приложения (API). Обработка данных происходит на стороне сервера, а сами данные и файлы хранятся в базе данных и файловой системе либо в специализированных хранилищах. Для реализации учебной видеоплатформы используется следующий стек технологий: фронтенд на React, бэкенд на Python, реляционная база данных PostgreSQL, а также Docker для упрощения развёртывания среды.

\textbf{Цель} расчётно-графической работы — разработка тестовой видеоплатформы, предоставляющей пользователю возможность просматривать список видеороликов, загружать собственные видео, изменять информацию о них и удалять записи. Также предусмотрена регистрация пользователя и последующее использование учётной записи при работе с приложением. Интерфейс должен быть удобным для пользователя, соответствовать базовым принципам UI/UX и корректно отображаться на устройствах с различной шириной экрана.

\textbf{Задачи работы:}
\begin{itemize}
    \item разработать и реализовать REST API на языке Python (с использованием фреймворка FastAPI), обеспечивающее хранение данных пользователей и информации о видео в базе PostgreSQL, а также загрузку и отдачу видеофайлов;
    \item реализовать механизм регистрации пользователя и хранение состояния авторизации на стороне клиента (с использованием React и \texttt{localStorage});
    \item создать клиентскую часть приложения на React с использованием маршрутизации (React Router), включающую каталог видео, страницу регистрации и страницу просмотра видео с комментариями;
    \item реализовать пользовательский интерфейс в соответствии с макетом, разработанным в Figma, с учётом цветовой схемы, расположения элементов и адаптивности;
    \item обеспечить удобную структуру проекта, подготовить инструкции по запуску и описать процесс разработки, включая возникшие трудности и способы их решения.
\end{itemize}

Практическая значимость работы заключается в закреплении навыков разработки современных веб-приложений, включая проектирование API, работу с базами данных, интеграцию клиентской и серверной частей, а также решение типичных проблем, возникающих при разработке (настройка CORS, работа с портами, взаимодействие с Docker и согласование адресов между фронтендом и бэкендом).